\documentclass[11pt,a4paper]{article}

\usepackage[T1]{fontenc}
\usepackage[utf8]{inputenc}
\usepackage{lmodern}
\usepackage[margin=2.6cm]{geometry}
\usepackage{microtype}
\usepackage{amsmath,amssymb}
\usepackage{booktabs}
\usepackage{array}
\usepackage{enumitem}
\usepackage{titlesec}
\usepackage[hidelinks]{hyperref}
\usepackage{xcolor}

\definecolor{rulegrey}{gray}{0.45}
\titleformat{\section}{\large\bfseries}{\thesection}{0.7em}{}
\titleformat{\subsection}{\normalsize\bfseries}{\thesubsection}{0.7em}{}
\titlespacing*{\section}{0pt}{1.6em}{0.6em}
\titlespacing*{\subsection}{0pt}{1.2em}{0.4em}
\setlist[itemize]{leftmargin=1.2em,itemsep=0.25em,topsep=0.3em}
\setlist[enumerate]{leftmargin=1.5em,itemsep=0.45em,topsep=0.4em}
\newcommand{\hi}[1]{\textbf{#1}}
\newcolumntype{d}{>{\raggedleft\arraybackslash}p{1.45cm}}

\begin{document}

\begin{center}
{\LARGE\bfseries The EM chapters, re-examined}\\[0.45em]
{\large Confirmatory evidence for Chapters 8--11}\\[0.9em]
{\small Internal report for discussion with Jerome and Marc \quad$\cdot$\quad 18 September 2026}\\[0.25em]
{\footnotesize Second revision. Changes against the 16 September draft are listed on this page.}
\end{center}

\vspace{0.4em}
\hrule height 0.8pt
\vspace{1.2em}

\noindent
This report states what sixteen-seed confirmatory runs, executed on the Bocconi cluster over the
past four days, say about the chapters from EM onwards. Three of the four chapters survive intact;
one conclusion reverses; one comparison the thesis explicitly declared out of reach has now been
made. Section~\ref{sec:summary} is the short version, Section~\ref{sec:method} states what we are
doing and why it is trustworthy, Sections~\ref{sec:ch8}--\ref{sec:ch11} go chapter by chapter,
Section~\ref{sec:robust} reports what four further runs say about how far the results travel, and
Section~\ref{sec:ask} lists the decisions that need taking.

\medskip
\noindent
\textbf{What changed since the 16 September draft.} Four runs have reported. Two of them confirm
and extend; two of them require corrections that are made here and not buried.
\begin{itemize}[topsep=0.25em]
\item The misspecification sweep was repeated at sixteen seeds. \textbf{One of its four
conclusions does not survive}: long-range coupling was reported as reversing the ordering near
$\gamma\approx0.4$; at sixteen seeds it does not reverse anywhere in the swept range. The
remaining three mechanisms are unchanged and two of them now have tighter brackets
(Section~\ref{sec:gen}).
\item The capacity contrast was re-run across five correlations and three chain lengths and holds
at all fifteen, with the apparatus checking itself at every one (Section~\ref{sec:regimes}).
\item A sweep of the quadrature found that the reference grid's accuracy is bounded by the
\emph{domain}, not by the number of points, and that one configuration the phase appendix reports
--- the heavy-tailed family at zero correlation --- lies outside the validated envelope
(Section~\ref{sec:grid}). Those entries are flagged and not cited pending a recomputation.
\item An audit of the EM iteration budget across all runs \textbf{flags one reported claim}: the
statement that the estimator leaves the information bound at the two highest noise levels is
confounded with the ladder running out there (Section~\ref{sec:budget}). The capacity conclusion
is cleared by the same audit.
\end{itemize}

\section{Summary}\label{sec:summary}

\begin{enumerate}
\item \hi{The capacity conclusion of Chapter~10 reverses.} It was an artefact of a stopping rule
that fired for the single-Gaussian arm and for no other. Under a protocol that selects the
checkpoint identically for every capacity, the single Gaussian sits exactly on the second-order
floor at every sample size ($0.167$) and every richer capacity goes far below it ($0.023$ at
$M=2048$), favouring the mixture in $16$ of $16$ seeds at every size.

\item \hi{The two rates of Chapter~9 are a property of the channel, not of the optimiser.} The EM
estimator attains the Cram\'er--Rao bound of the noised likelihood to within a factor $1.02$ on
average. The shape coordinate is not slow because EM is slow; the data stop carrying it. Its
information bound crosses the coordinate's own range at $t=0.98$, and that is precisely where the
fit collapses to a Gaussian innovation.

\item \hi{The comparison of Chapter~10 now runs at a matched information interface.} The thesis
attached a scope condition saying it could not do this. Given the same single noised view per
sequence that the estimator sees, the best network reaches $0.155$ against the estimator's
$0.023$: a paired factor $7.3$, interval $[6.2,8.7]$. The unmatched comparison, at $2.83\times$,
\emph{understates} the gap.

\item \hi{Locality is a ceiling, not only a prior.} The radius-4 window head is the best network at
$M=32$ and the worst of the five weight-shared architectures at $M=2048$, moving only
$0.121\rightarrow0.118$ as data quadruples. A bounded receptive field cannot represent a message
that has crossed the chain --- which is exactly what the chain's score is.

\item \hi{Score accuracy transfers to generation, and the second-order score fails there
visibly.} Driven through the same reverse SDE on identical Brownian paths, the fitted kernel
produces sequences whose innovation excess kurtosis is $1.887$ against the reference's $1.896$,
with a log-likelihood gap of $0.022$ under the true law. The covariance-matched score gives
$-0.008$ --- an essentially Gaussian chain --- at a gap of $1.71$, and the isotropic score a gap of
$46.2$.

\item \hi{A claim of ours was refuted by our own data, and is recorded as such.} ``The learned
networks sit at or above the second-order floor'' had been drafted into the new abstract. At
$M=2048$ five of six architectures sit below it. The defensible statement is about the information
interface, not about a barrier the networks cannot pass.

\item \hi{The capacity result travels, and checks the apparatus while it does.} Across five
correlations and three chain lengths the mixture beats the Gaussian kernel in all fifteen cells,
with every interval excluding zero and the margin shrinking smoothly with $\rho$ as the marginal's
excess kurtosis does. In the same fifteen cells the $C=1$ arm lands on the LMMSE of the true
covariance to $0.01$, as Proposition~1 requires --- a simultaneous check on the grid, the message
passing, the EM and the scoring, passed at fifteen independent configurations.

\item \hi{Two corrections, both found by us and both made here.} Long-range coupling does
\emph{not} reverse the ordering at sixteen seeds, though the eight-seed draft said it did; and
the statement that the estimator leaves the information bound at the two highest noise levels is
confounded with the EM ladder running out there. The first is settled and rewritten, the second is
flagged with a control queued. A third finding of the same kind is about our instrument rather
than our results: the reference grid's accuracy is limited by the width of its domain rather than
by the number of points, and one configuration the phase appendix reports lies outside the
validated envelope.
\end{enumerate}

\section{What we are doing, and why the numbers can be trusted}\label{sec:method}

Three changes to how the work is run account for why these results settle questions the thesis
left open.

\paragraph{One frozen protocol, one varied factor.}
Everything the arms share is now a single frozen object: the bundle sizes and their random
streams, the sizes $M\in\{32,128,512,2048\}$, the sixteen confirmatory seeds and four development
seeds, the EM iteration ladder and its initialisation, the network training recipe and screen
budget, the estimand, and the scoring regions. Each experiment varies exactly one declared factor
against that background, so a difference between two arms is attributable to that factor. Two arms
that differ in more than one thing are not compared; where a comparison genuinely needs two
factors to move, the intermediate arm is run as well.

\paragraph{Pairing by construction.}
Every row records the SHA-256 of the validation and test bundles it was scored on. Rows from
different experiments are joined on (seed, size, bundle hash), and a generator that is handed rows
with mismatched hashes refuses to emit a number. This is what makes ``EM--BP against the best
network'' a paired comparison across two experiments rather than two averages.

\paragraph{Frozen outputs.}
A result becomes citable only when its output directory is pinned: every CSV and provenance file
hashed into the registry entry together with the commit that produced it. After that, any byte
change fails the build, and the manuscript refuses to compile if a generated number disagrees with
the artefact it came from. A re-run is a new artefact under a new key, never an overwrite.
Confirmatory seeds are declared before the run and never used for selection; development seeds
never enter a reported number.

\paragraph{Amendments are recorded, not absorbed.}
Two design changes were made this week before the runs that used them, and both are in the
registry with timestamps: the EM ladder was extended from $400$ to $1200$ iterations because $400$
was binding at the largest sample size, and the largest strength of the site-varying
misspecification mechanism was removed because it is not a valid chain --- at $\rho=0.85$ its site
coefficients exceed $1$ and the implied covariance is not positive definite. No number from the
invalid point exists.

\section{Chapter 8 --- EM for the parameters: holds}\label{sec:ch8}

Nothing in the derivation needed revision. Fisher's identity still gives the exact gradient of the
discretised marginal log-likelihood from one forward--backward pass, and the finite-difference
check agrees to about $10^{-9}$. Two things were added.

The algebra is now machine-checked where checking is cheap. A Lean~4 module builds against Mathlib
with no \texttt{sorry} and no axiom, covering four constants: the Hermite attenuation for degrees
up to four, the fourth-cumulant difference under variance matching, the factor $\tfrac12$ in the
Gaussian information rate, and the coefficient $\kappa_4/6$ in the score gap. It verifies the
algebraic skeleton and not the analysis: the remainder bounds and the measure-level statements
remain informal, and the write-up says so in those words rather than claiming more.

The confirmatory ladder was extended to $1200$ iterations. At $M=2048$ the richer capacities were
still selecting the top of the $400$-iteration ladder; they still select the top of the extended
one in $10$ of $16$ seeds, so the errors reported for those arms in Table~\ref{tab:capacity}
remain upper bounds. That can only understate the effect described there.

\section{Chapter 9 --- the two rates: same fact, better reason}\label{sec:ch9}

The chapter explains the asymmetry through EM's fixed-point map: the Jacobian is the missing
information and the shape coordinate has the smaller eigenvalue, so the fitted innovation shape
settles at a median of $229$ updates against $80$ for the correlation. That measurement stands,
and the stopping-rule warning built on it stands with it.

What is new is that the asymmetry is no longer only a statement about the optimiser. Two results
bracket it.

\paragraph{The channel, from theory.}
Under the Ornstein--Uhlenbeck channel every Hermite mode of degree $m$ is attenuated exactly by
$e^{-mt}$, so the Fisher information one noised sequence carries about order-$m$ structure decays
as $e^{-2mt}$: correlation as $e^{-4t}$, non-Gaussian shape as $e^{-8t}$. Computed on the chain by
exact inference the measured slopes are $-3.97$ and $-8.01$ against the predicted $-4$ and $-8$,
and the constant in the score gap between the true score and that of the covariance-matched
Gaussian is reproduced to within one per cent.

\paragraph{The estimator, from measurement.}
Fitting the transition law from sequences observed at a single level $t$ ($M=512$, $16$ seeds, one
noised view per sequence), the EM estimator sits \emph{at} the Cram\'er--Rao bound of the noised
likelihood: a ratio of $1.02$ on average over the levels where either coordinate is informative,
never worse than $1.18$ for the correlation or $1.51$ for the fourth cumulant.

\medskip
\noindent
Taken together these say the chapter's phenomenon belongs to the channel. If the estimator is
efficient, the two rates cannot be blamed on EM. The bound for the shape coordinate grows as
$e^{4t}$ and crosses the coordinate's own range ($\kappa_4=3$) at $t=0.98$; the fit collapses to a
Gaussian innovation at exactly that level, while the correlation estimate is still within $0.009$
of the truth and survives past the crossover $t_*=1.17$.

\paragraph{One limitation, stated plainly.}
The corollary predicts errors growing as $e^{2t}$ and $e^{4t}$, and this run cannot test those
exponents. At $\rho=0.85$ the shape estimate saturates before the crossover, so the asymptotic
window does not exist at this configuration. Fitting slopes through a saturated estimate would
have returned $3.64$ and $0.09$ against predictions of $2$ and $4$; the analysis therefore reports
no slope, and a companion run at $\rho=0.5$ with $M=2048$ --- where the crossover falls at $0.55$
and the window is reachable --- is registered and queued.

\section{Chapter 10 --- the comparison: one reversal, one replacement, two new facts}\label{sec:ch10}

\subsection{The capacity sweep reverses}

The chapter concluded that no mixture capacity improves on a single Gaussian innovation, and
flagged its own caveat in the accompanying figure: the share of fits stopped at the iteration cap,
and with a narrowest component below the grid's resolution floor, both climbed steeply with $C$.
That caveat was the entire effect. In the archival sweep the likelihood tolerance fired in $31$ of
$32$ single-Gaussian cells at a median of $71$ iterations --- close to where a validation protocol
would stop --- and in none of the $96$ cells with $C\ge4$, which therefore reported checkpoints an
order of magnitude past their optimum. ``The single Gaussian is as good'' was ``the single Gaussian
is the only arm that stopped in the right place''.

Table~\ref{tab:capacity} re-runs the sweep with the checkpoint selected on a per-seed validation
bundle identically for every capacity.

\begin{table}[htbp]
\centering
\small
\caption{Capacity under equal stopping. Schedule-pooled relative score error on all sites, mean
over $16$ seeds, every arm on identical held-out bundles. The last column is the same capacity
fitted to the innovation law itself --- the $C$-component mixture closest in Kullback--Leibler
divergence --- an infinite-data reference for the family rather than a bound. The isotropic score
is at $0.669$ at every size.}
\label{tab:capacity}
\begin{tabular}{l dddd c d}
\toprule
& \multicolumn{4}{c}{sequences $M$} & & \\
\cmidrule(lr){2-5}
kernel & 32 & 128 & 512 & 2048 & & $M=\infty$ \\
\midrule
LMMSE floor      & 0.167 & 0.167 & 0.167 & 0.167 & & 0.167 \\
$C=1$            & 0.174 & 0.168 & 0.168 & 0.167 & & 0.167 \\
$C=2$            & 0.084 & 0.050 & 0.035 & 0.028 & & 0.039 \\
$C=4$            & 0.093 & 0.053 & 0.037 & 0.025 & & 0.0057 \\
$C=8$            & 0.099 & 0.055 & 0.036 & 0.023 & & 0.0022 \\
$C=16$           & 0.100 & 0.056 & 0.036 & 0.022 & & 0.0007 \\
\bottomrule
\end{tabular}
\end{table}

Three readings follow. The single-Gaussian kernel lands on the second-order floor at every sample
size, which is Proposition~1 read as a statement about an estimator rather than about a
projection. Every richer capacity goes below that floor, by a factor $7.9$ at the largest size; the
paired difference between $C=1$ and $C=8$ is $0.144$ with bootstrap interval $[0.141,0.147]$ and
favours the mixture in $16$ of $16$ seeds. And the best capacity rises with the data, from $C=2$ at
$M=32$ to $C=16$ at $M=2048$ --- the ordinary variance-to-bias trade, visible now that the stopping
rule no longer masks it.

\paragraph{The check that closes the loop.}
Two things changed at once between the thesis sweep and this one: the protocol, and the estimand
(held-out log-evidence per edge there, relative score error here). Two simultaneous changes are not
an explanation. A registered run now scores the \emph{same frozen fits} in the thesis's own
estimand, leaving the protocol as the only difference. A single-seed pass already shows the
reversal --- log-evidence per edge $-1.3116$ for $C=8$ against $-1.3146$ for $C=1$, with the true
kernel at $-1.3076$ --- and the sixteen-seed result is running.

\subsection{The data-efficiency comparison is replaced}

The chapter reports $6.6$--$10.9\times$ against unstructured denoising score matching and
$2.3$--$6.1\times$ against a baseline given locality and weight sharing, with three scope
conditions, the sharpest being that the arms are asymmetric in supplied information and that the
comparison therefore does not rank learning algorithms at a matched interface.

Six architectures were screened under one recipe --- $68$ configurations, four development seeds,
three restarts, $32{,}000$ steps each --- and the winner of each family was run on the sixteen
confirmatory seeds (Table~\ref{tab:baselines}).

\begin{table}[htbp]
\centering
\small
\caption{Under one training recipe, $16$ confirmatory seeds, identical held-out bundles. All-site
schedule-pooled relative score error. The upper block is trained on clean sequences with a fresh
corruption at every gradient step; the lower block sees what the estimator sees.}
\label{tab:baselines}
\begin{tabular}{l dd l}
\toprule
estimator & $M=32$ & $M=2048$ & information supplied \\
\midrule
fully connected              & 0.517 & 0.208 & clean, fresh corruption \\
window head ($r=4$)          & 0.186 & 0.118 & clean, fresh corruption \\
bidirectional GRU            & 0.201 & 0.060 & clean, fresh corruption \\
transformer                  & 0.231 & 0.063 & clean, fresh corruption \\
1-D U-Net                    & 0.227 & 0.066 & clean, fresh corruption \\
dilated TCN                  & 0.223 & 0.068 & clean, fresh corruption \\
\midrule
best matched-interface net   & ---   & 0.155 & one noised view \\
EM--BP ($C=8$)               & 0.099 & 0.023 & one noised view \\
\midrule
\textit{LMMSE floor}         & \textit{0.167} & \textit{0.167} & --- \\
\bottomrule
\end{tabular}
\end{table}

Two numbers replace the thesis's ranges. Against the best of six architectures under the standard
recipe, the estimator is ahead by a paired factor $2.83$, interval $[2.37,3.42]$, in $16$ of $16$
seeds at $M=2048$, and by $2.07$ at $M=32$. Against the best network trained at the estimator's own
interface --- one noised view per sequence, no clean chain, checkpoint selected on the implicit
objective of a held-out noisy bundle and never on the true score --- the factor is $7.3$, interval
$[6.2,8.7]$. Implicit score matching on that bundle, which uses no clean target at all, lands at
$0.165$.

The second number is the one the thesis said was out of reach, and it is the stronger claim: the
gap is not an artefact of giving the networks less, because when they are given the same, the gap
widens.

It also follows that one sentence drafted for the new paper is false and has been withdrawn.
``Every network sits at or above the second-order floor'' holds only at the smallest sample size;
at $M=2048$ five of the six sit below it. Given enough clean sequences and fresh corruptions a
neural denoiser plainly does learn fourth-order structure. The honest statement concerns the
information interface.

\subsection{Two facts the thesis does not contain}

\paragraph{Locality is a ceiling.}
The window head is the best network at $M=32$ and the worst of the five weight-shared
architectures at $M=2048$, moving only from $0.121$ to $0.118$ as the data quadruples, against a
standard error of $0.001$, while the global architectures continue to $0.060$. A radius-4 window
cannot represent a message that has travelled the length of the chain, and the chain's exact score
is precisely such a message. This connects the architecture comparison to the bulk-posterior
analysis the chapter already contains.

\paragraph{Length extrapolation separates the arms cleanly.}
Fitted at $L=32$ and evaluated against the exact score of longer chains
(Table~\ref{tab:length}), the fitted kernel is flat: the law is local, so inference at any length
is the same exact recursion. The convolutional and recurrent networks improve slightly, as they
should, since longer chains have proportionally fewer boundary sites. The transformer with learned
relative positions degrades by an order of magnitude --- a failure of the positional encoding
beyond its training length, not of the score it learned.

\begin{table}[htbp]
\centering
\small
\caption{Length extrapolation at $M=2048$: pooled relative score error of estimators fitted at
$L=32$ and evaluated against the tree-exact score at longer chains, mean over $16$ seeds.}
\label{tab:length}
\begin{tabular}{l dddd}
\toprule
estimator & $L=32$ & $L=64$ & $L=128$ & $L=256$ \\
\midrule
EM--BP ($C=8$)     & 0.023 & 0.023 & 0.023 & 0.023 \\
bidirectional GRU  & 0.060 & 0.057 & 0.057 & 0.056 \\
1-D U-Net          & 0.066 & 0.064 & 0.062 & 0.062 \\
dilated TCN        & 0.068 & 0.066 & 0.065 & 0.064 \\
window head        & 0.118 & 0.099 & 0.088 & 0.082 \\
transformer        & 0.063 & 0.243 & 0.514 & 0.703 \\
\bottomrule
\end{tabular}
\end{table}

\section{Chapter 11 --- which answers change}\label{sec:ch11}

\paragraph{RQ5, recovering an unknown transition.}
The answer holds and gains a sentence. The E-step is structurally exact, recovery is demonstrated
numerically over the studied regime, and the uneven settling of the coordinates is real. We can now
add that the estimator is statistically efficient, so the unevenness belongs to the channel and not
to the algorithm. The stopping-rule warning becomes sharper rather than weaker: a rule read off the
correlation trace certifies a kernel whose shape coordinate the data may no longer identify at all.

\paragraph{RQ6, what supplying the Markov assumption is worth.}
The answer changes in its numbers and in its scope. The ranges $6.6$--$10.9\times$ and
$2.3$--$6.1\times$ are superseded by $2.83\times$ against the best of six architectures under one
recipe. The third scope condition --- that the comparison does not rank methods at a matched
information interface --- can be removed, because that comparison now exists and gives
$7.3\times$. The capacity subsection of the same chapter must be rewritten in the opposite
direction.

\paragraph{The retired claim.}
The capacity result should be stated as retired rather than quietly dropped, together with the
diagnosis of why it arose. The same applies to the claim of ours that our own screen refuted this
week. The project keeps a negative-results record for exactly this, and both entries are in it.

\section{Generation and the boundary of the advantage}\label{sec:gen}

Two results completed while this report was being written. Both were pre-registered and both
address questions Chapter~10 raises but does not settle.

\paragraph{Does score accuracy transfer to generation?}
Every arm drives the same reverse SDE from the same initialisation along the same Brownian path, so
the comparison is paired within a seed and the integrator is common to all of them. The
discretisation was measured separately and is not what separates the arms: against an $800$-step
reference the strong error is $0.0190$, $0.0087$ and $0.0033$ at $100$, $200$ and $400$ steps,
roughly halving per doubling, while the excess kurtosis of the generated sequences is stable at
$2.00$--$2.05$ across the whole ladder.

\begin{table}[htbp]
\centering
\small
\caption{Generation at $M=2048$, sixteen seeds, every arm on the same Brownian path within a seed.
The deficit is the paired difference in mean $\log p_{t_{\min}}$ per sequence against the arm driven
by the exact score, in nats, with a seed bootstrap; the kurtosis column is $|$generated $-$
reference$|$, the reference sample having excess kurtosis $1.936$.}
\label{tab:generation}
\begin{tabular}{l l dd}
\toprule
score driving the sampler & deficit vs exact (nats) & \multicolumn{1}{c}{kurt.\ error} & \multicolumn{1}{c}{cov.\ error} \\
\midrule
exact (tree BP)         & ---                        & 0.119 & 0.070 \\
EM--BP ($C=8$), fitted  & $+0.09$\ \ $[+0.00,+0.19]$ & 0.177 & 0.077 \\
bidirectional GRU       & $+0.32$\ \ $[+0.17,+0.46]$ & 0.385 & 0.102 \\
window head             & $+1.05$\ \ $[+0.79,+1.31]$ & 0.474 & 0.290 \\
LMMSE (covariance twin) & $+1.70$\ \ $[+1.68,+1.71]$ & 1.939 & 0.069 \\
isotropic               & $+46.2$                    & 1.936 & 0.908 \\
\bottomrule
\end{tabular}
\end{table}

The ordering is the one the score errors predict, and the fitted kernel is nearest the exact score
by a wide margin: its deficit is $0.09$ nats against $0.32$ for the best network and $1.70$ for the
covariance twin. The deficit is small but real --- an earlier four-seed run suggested the fitted
score matched the exact one, and sixteen seeds show it does not. The second-order score is the
instructive failure. On the covariance error it is the best arm in the
table, marginally ahead of the exact score, which is what Proposition~1 predicts. On the fourth
cumulant it is the worst: it generates an essentially Gaussian chain with the right covariance. A
generative model evaluated only on second-order distances would score it a success.

\paragraph{Where does the advantage stop?}
The structured estimator assumes a first-order Markov chain. Four Gaussian departures from that
assumption were swept, each with a scalar strength $\varepsilon$ and an exact score of its own, so
every arm is scored against the truth of the perturbed model rather than against the chain.

\begin{table}[htbp]
\centering
\small
\caption{Misspecification at sixteen seeds: ratio of the window head's pooled error to EM--BP's,
at zero strength and at the largest strength swept, $M=2048$. Above $1$ the misspecified chain
estimator still wins. The break-even is reported as the bracket between the two grid points that
straddle it, with the interpolated crossing in parentheses; we do not fit a crossing. The
long-range row is the correction against the eight-seed draft.}
\label{tab:misspec}
\begin{tabular}{@{}l l d d d l@{}}
\toprule
mechanism & what it breaks & \multicolumn{1}{c}{$\varepsilon=0$} & \multicolumn{1}{c}{$\varepsilon_{\max}$} & \multicolumn{1}{c}{ratio} & break-even \\
\midrule
rank-one global latent   & nothing structural & 15.55 & 1.00 & 3.97 & not reached \\
AR(2) second coefficient & the order          & 20.78 & 0.40 & 0.96 & $(0.3,\,0.4)$, $0.38$ \\
long-range coupling      & the precision band & 14.11 & 0.40 & 1.00 & not reached \\
site-varying $\rho$      & homogeneity        & 14.22 & 0.14 & 0.98 & $(0.10,\,0.14)$, $0.14$ \\
\bottomrule
\end{tabular}
\end{table}

The pattern is interpretable, and it is not quite the one the eight-seed draft reported. A
rank-one global latent is absorbed by the fitted correlation and never reverses the ordering,
because it does not break the conditional independence the estimator exploits; that stands. Two
mechanisms that do break it --- adding an order, and making the chain inhomogeneous --- reverse it
at a measurable strength, and the brackets are now tight.

Long-range coupling is the correction. At eight seeds its ratio fell below $1$ at
$\gamma=0.4$ and was reported as a reversal near that point. At sixteen it does not: the paired
log-ratio is $+0.004$ with interval $[-0.000,\,0.009]$, a tie rather than a crossing, at both
sample sizes. The eight-seed reading was inside its own noise. This also retires an inherited
claim: the thesis reports the advantage reversing under long-range coupling near $\gamma\approx
0.1$, from a single draw per cell against an unscreened baseline. Against a screened baseline with
sixteen seeds there is no reversal anywhere up to $\gamma=0.4$.

What survives is the boundary the chapter lacked, stated more carefully: the advantage is not a
general property of ``having structure'', it is specific to the structure being approximately
right; breaking the order or the homogeneity costs it at a strength we can now name, and breaking
the precision band or adding a global latent does not cost it within the range we swept. The two
sample sizes also separate the two failure modes as predicted --- at zero strength the EM--BP error
halves from $0.019$ to $0.010$ between $M=512$ and $M=2048$, which is estimation; at the largest
AR(2) strength it moves from $0.2135$ to $0.2136$, which is bias, and no amount of data will move
it.

\section{How far the results travel}\label{sec:robust}

Everything above is measured at one point of the model's parameter space: a Laplace innovation,
$\rho=0.85$, chain length $32$. Four further runs ask whether that matters. Two of them strengthen
the chapters; two of them require corrections, which are made here.

\subsection{The capacity contrast holds in every regime}\label{sec:regimes}

The capacity result is the load-bearing one, so it was re-run across five correlations
$\rho\in\{0.3,0.5,0.7,0.85,0.95\}$ and three chain lengths $L\in\{16,32,64\}$, eight seeds each,
everything else at the frozen protocol.

\begin{table}[htbp]
\centering
\small
\caption{Capacity across regimes: mean pooled relative score error of the Gaussian-innovation
kernel ($C=1$) and the mixture ($C=8$), and their paired difference with a percentile bootstrap
over seeds. Negative favours the mixture. Every interval excludes zero. $L=64$ shown; the shorter
chains agree.}
\label{tab:regimes}
\begin{tabular}{l d d d l d}
\toprule
$\rho$ & \multicolumn{1}{c}{$t_*$} & \multicolumn{1}{c}{$C=1$} & \multicolumn{1}{c}{$C=8$} & paired difference & \multicolumn{1}{c}{$\kappa_4$ of marginal} \\
\midrule
0.30 & 0.31 & 0.286 & 0.033 & $-0.253$\; $[-0.259,-0.246]$ & 2.50 \\
0.50 & 0.55 & 0.266 & 0.035 & $-0.231$\; $[-0.239,-0.224]$ & 1.80 \\
0.70 & 0.86 & 0.227 & 0.032 & $-0.195$\; $[-0.198,-0.192]$ & 1.03 \\
0.85 & 1.22 & 0.167 & 0.028 & $-0.139$\; $[-0.142,-0.135]$ & 0.48 \\
0.95 & 1.67 & 0.086 & 0.025 & $-0.061$\; $[-0.064,-0.058]$ & 0.15 \\
\bottomrule
\end{tabular}
\end{table}

Two things are worth having. First, the result is a property of the estimator and not of the
corner: the mixture wins in all fifteen cells with the interval excluding zero, and the margin
moves smoothly with $\rho$ rather than appearing at one point. Second, and more useful to us, the
sweep audits the apparatus at every cell. Proposition~1 says the $C=1$ arm \emph{must} coincide
with the LMMSE of the true covariance; it does, to $0.002$ in the median cell and $0.010$ in the
worst, and its fitted excess kurtosis is zero to three decimals everywhere. That is a check on the
grid, the message passing, the EM and the scoring simultaneously, repeated at fifteen independent
configurations, and it passes at all of them. The $C=8$ arm recovers the innovation's true excess
kurtosis of $3$ at every correlation ($2.85$ to $3.40$), with its spread across seeds widening from
$\pm0.33$ at $\rho=0.3$ to $\pm1.38$ at $\rho=0.85$ --- the shape gets harder to see, not
differently placed, which is the same information story from a third direction.

The margin shrinks monotonically with $\rho$, in the direction the theory gives: a strongly
correlated chain has a nearly Gaussian marginal, its excess kurtosis falling as
$3(1-\rho^2)/(1+\rho^2)$ (last column), so there is less fourth-order structure for capacity to
buy. We report that as a direction and not as a rate. The pooled error runs over a schedule whose
informative window itself moves with $\rho$, so its trend is not the exponent of the closure
theorem, and fitting one to it would be over-reading.

\subsection{The reference grid is not uniformly converged}\label{sec:grid}

Every number in this project is measured against tree-exact belief propagation on a grid of $401$
points on $[-8,8]$. That choice was validated once, at the frozen configuration; every other
configuration we touch has inherited the validation without earning it. A sweep of five innovation
families $\times$ four correlations $\times$ twelve levels $\times$ seven candidate grids against a
$(3201,\,12)$ reference now measures it, and the result has a reassuring half and an uncomfortable
half.

The reassuring half: \emph{at the frozen configuration the operating grid is excellent}. The
pooled drift of the posterior-mean field against the reference is $9.2\times10^{-6}$, the worst
single level $3.7\times10^{-5}$, and $2.8\times10^{-4}$ of the belief mass sits in the outermost
$2\%$ of the domain. Nothing in Sections~\ref{sec:ch8}--\ref{sec:ch11} is in question.

The uncomfortable half: over the whole sweep, the worst-case drift is \emph{flat in the number of
grid points}. At half-width $8$ it is $5.411\times10^{-2}$ at $201$ points and
$5.425\times10^{-2}$ at $1601$ --- refining by a factor of eight buys nothing --- and it falls
only when the domain widens, to $2.6\times10^{-2}$ at $(401,\,10)$ and $3.0\times10^{-3}$ at
$(801,\,12)$. That is the signature of truncation, not of quadrature. It matters because
``refinement'' in this project, and in the thesis before it, has meant increasing the number of
points, and that operation cannot fix this error.

The binding cell is Student-$t_5$ at $\rho=0$: essentially all of its posterior mass falls outside
$[-8,8]$, and the drift is $5.4\times10^{-2}$. \emph{That cell is reported}, in the phase appendix
that orders the families by kurtosis at each correlation. Its entries are now flagged in the
manuscript, marked in the figure, and excluded from the ordering claim; two recomputations at
$(801,\,12)$ and $(1601,\,16)$ are queued to establish what they should be. No result in the
chapters above depends on it --- the chapters live at $\rho=0.85$ with a Laplace innovation --- but
it is a reported number outside its validated envelope, and it should not have been.

\subsection{The iteration budget, audited across every run}\label{sec:budget}

The capacity reversal was caused by a stopping rule. That made it worth asking, of every run we
have, how often the validation-selected checkpoint is the \emph{last} rung of the EM ladder ---
which is the signal that validation was still improving when the ladder ran out, and that a longer
ladder might have selected differently.

\begin{table}[htbp]
\centering
\small
\caption{Fraction of cells whose validation-selected checkpoint is the last rung of the
$1200$-iteration ladder. A high rate does not by itself invalidate a result; it means the result is
a statement about a budget until something shows otherwise.}
\label{tab:atcap}
\begin{tabular}{@{}l l >{\raggedright\arraybackslash}p{6.1cm}@{}}
\toprule
run & rate & reading \\
\midrule
E02a (capacity)         & up to $0.69$ at $M=2048$            & \emph{cleared}: the $400$-iteration replication reproduces every overlapping cell to $0.001$, so the estimates have plateaued well before the cap \\[0.35em]
E13 (regimes)           & $\le 0.25$; $C=1$ never             & no action \\[0.2em]
E08 (misspecification)  & $\approx 0.10$                      & no action \\[0.2em]
E06 (learning)          & $0.12$--$0.38$ for $t\le1$          & \emph{cleared} at the informative levels \\
                        & $0.88$ at $t=1.4$, $0.75$ at $t=2$  & \textbf{flagged}, see below \\[0.2em]
E12 (sample complexity) & $1.00$ at $t=1$, $M\ge4096$         & \textbf{confounded}, see below \\
\bottomrule
\end{tabular}
\end{table}

\paragraph{The flagged claim.} Section~\ref{sec:ch9} states that the estimator attains the
Cram\'er--Rao bound over the informative levels and leaves it past the crossover, at $t=1.4$ and
$t=2$. The first half is safe: at every level up to $t=1$, validation stops of its own accord well
inside the ladder, so the budget is not binding and the efficiency result stands. The second half
is not: at exactly the two levels where we report the estimator leaving the bound, it is also
running out of iterations. The departure may be information or it may be budget, and this run
cannot tell which. The manuscript now says so, and a control with the ladder extended sixfold is
queued. Nothing else in the chapter depends on it.

\paragraph{The confounded run.} The sample-complexity experiment was designed to escape the
saturation problem by asking how much data buys back a fixed accuracy as the level rises --- a
quantity that does not saturate. At $t=0.6$ and $t=0.8$ it worked. At $t=1.0$ it did not: the
error is flat and \emph{non-monotone} in $M$, and the fitted kurtosis scatters between $0$ and
$3.4$ across seeds with no concentration, because the EM is censored by the iteration cap in forty
cells out of forty. At that level the experiment measured the optimiser's budget, not the sample
size.

Two one-sided statements survive, and they are worth having because they come from opposite
directions: the covariance coordinate's log-slope is $5.40$ with interval $[2.13,\,7.57]$, which
\emph{excludes} the shape prediction of $8$ and contains its own prediction of $4$; and the shape
coordinate fails to reach the accuracy threshold at $t=1$ even at $M=8192$, which forces its
log-slope above $5.54$ and \emph{excludes} $4$. The two coordinates are therefore separated, which
is what the corollary asserts. The point estimates where a crossing exists sit on the predictions
--- $7.54$ and $8.53$ against $8$, and $4.44$ for the paired ratio against $4$ --- but their
bootstrap intervals exclude nothing at all, and we do not offer them as evidence. A re-run with
five levels, sixteen seeds and a ladder extended to $8000$ is queued.

\section{Status of the evidence, and what is still running}\label{sec:status}

\begin{table}[htbp]
\centering
\small
\caption{Confirmatory runs. \emph{Frozen} means every file in the output directory is pinned by
SHA-256 in its registry entry, with the commit that produced it; a generator refuses to read an
output that is not frozen or whose bytes have changed.}
\label{tab:status}
\begin{tabular}{l l l}
\toprule
run & what it establishes & state \\
\midrule
E01  & order-by-order decay verified on the chain              & frozen \\
E02a & capacity under equal stopping                            & frozen \\
E02b & capacity against innovation shape at fixed covariance    & frozen \\
E03  & six architectures, matched interface, length             & frozen \\
E04  & grid refinement on the fitted models                     & frozen \\
E05  & generation under each score, paired noise                & frozen \\
E06  & learning through noise, against the information bound    & frozen \\
E07  & the second-order floor across correlation                & frozen, one cell flagged \\
E08  & misspecification and its break-even (16 seeds)           & frozen \\
E09  & infinite-data reference per capacity                     & frozen \\
E10  & capacity in the thesis's own estimand                    & frozen \\
E13  & the capacity contrast across $\rho$ and $L$              & frozen \\
\midrule
E06b & the asymptotic exponents at $\rho=0.5$                   & frozen, inconclusive \\
E11  & the quadrature envelope across families and $\rho$       & frozen; found the exclusion \\
E12  & sample complexity by order                               & frozen, confounded \\
\midrule
E06c & the same nine levels, ladder extended sixfold            & queued \\
E07b, E07c & the family ladder at $(801,12)$ and $(1601,16)$    & queued \\
E11b & the quadrature envelope at $L=64$ and $128$              & queued \\
E12b & five levels, sixteen seeds, ladder to $8000$             & queued \\
E14  & the matched comparison across $\rho$ and $L$             & queued \\
\bottomrule
\end{tabular}
\end{table}

\emph{Refinement} is complete and is as clean as this kind of check gets. Over the whole ladder ---
$401$, $801$ and $1601$ grid points, half-widths $8$ and $10$ --- every arm's error agrees to five
decimal places and the reference field itself moves by at most $10^{-5}$. No number in this report
is an artefact of the quadrature at the frozen configuration. That sweep has since been extended to
the other families and correlations, where the answer is less comfortable and is reported in
Section~\ref{sec:grid}.

Seven jobs are queued on the cluster against pre-registrations written before submission. Four are
controls for the two problems above; one carries the quadrature envelope to the longer chains. The
sixth and seventh are the one addition that is about strength rather than repair: \textbf{the
information-matched comparison, re-run across three correlations and two chain lengths with the
networks re-screened inside each regime}. That comparison is currently stated at a single point,
and it is the claim a referee will press on hardest; the capacity result has now been shown to
travel, and the comparison against networks has not.

\paragraph{Three infrastructure faults, for the record.}
Jobs were inheriting the login shell's $600$-second CPU limit, which killed one run outright and
left another holding an idle compute node for ten hours while the scheduler reported it running.
Two drivers wrote their results only at the very end, so a timeout would have lost everything; both
now write per cell and resume. And the site-varying misspecification mechanism is undefined at its
largest registered strength. All three are fixed; the third is a registered amendment.

\section{Decisions requested}\label{sec:ask}

\begin{enumerate}
\item \hi{How should the capacity reversal be presented?} Either rewrite the capacity section of
Chapter~10 around the new result and retain the old figure as a documented artefact, or keep both
and make the diagnosis itself the finding. The second is more honest about how the error surfaced
and costs roughly a page.

\item \hi{Is the matched-interface comparison the headline, or a robustness check?} It is the
stronger claim and it removes a scope condition the thesis carried. Promoting it changes the
framing from ``structure beats networks'' to ``structure beats networks at equal information, and
the unmatched comparison understates it''.

\item \hi{Where does the efficiency result belong?} The estimator attaining the Cram\'er--Rao bound
is what turns the two rates from an observation about EM into a statement about information. It
reads naturally in Chapter~9 or in the theory chapter, and should be stated once with a reference
from the other.

\item \hi{How far should the Lean verification go?} Four algebraic constants are machine-checked.
The remainder bounds are analysis and would be a substantial effort; the projection operator behind
Proposition~1 is the one load-bearing step still in prose. Is the current scope the right stopping
point?

\item \hi{Do we pursue the asymptotic exponents, or claim the rates only where they are
established?} Two attempts have now failed for two different reasons. The error-at-fixed-$M$ route
saturates: the window at $\rho=0.5$ is three levels wide, too narrow to pin an exponent
(Section~\ref{sec:ch9}). The sample-complexity route avoids saturation by construction but was
confounded by the iteration budget at its highest level (Section~\ref{sec:budget}); the re-run
that removes the confound is queued, and if it lands we will have the exponents, and if it does
not we will have two one-sided statements that separate the coordinates without naming a rate. The
fallback, which is what the write-up currently does, is to establish the rates exactly by
computation and the estimator's efficiency by measurement, and to state the sample-complexity
consequence only where the coordinate is identifiable. \emph{Our recommendation: hold the claim at
the fallback and let the re-run upgrade it if it can, rather than write the stronger version and
retract it.}
\item \hi{Does the misspecification boundary belong in Chapter~10 or in the conclusions?} It is
the cleanest statement we have about when the method should \emph{not} be used --- breaking the
order or the homogeneity costs the advantage at a strength we can name, while a global latent and
a widened precision band do not cost it within the range swept. It reads as a limitation, but it
is a measured one, which is unusual enough to be worth foregrounding.

\item \hi{Do we wait for the regime sweep of the matched comparison before fixing the framing?}
The headline factor of $7.3$ is stated at one correlation and one chain length. The capacity
result has now been shown to travel; this one has not been tested that way, and it is the claim a
referee will press on. The run is queued and needs roughly two days of cluster time.
\emph{Our recommendation: yes} --- the cost is small against the risk of building a framing on a
single point, and the sweep can only strengthen or usefully qualify it.

\item \hi{How should the quadrature exclusion be handled in the write-up?} One reported cell lies
outside the validated envelope (Section~\ref{sec:grid}). Three options: recompute and replace it
silently once the queued runs land; recompute and report the movement as a numerical appendix; or
keep the exclusion visible as a finding about the instrument, since ``refining the grid does not
fix a domain error'' is a real and transferable observation about this class of method.
\emph{Our recommendation: the second} --- honest, short, and it does not inflate a housekeeping
fix into a contribution.
\end{enumerate}

\vspace{1.4em}
\hrule height 0.4pt
\vspace{0.7em}
\noindent
{\footnotesize Every number above comes from a registered experiment whose output directory is
pinned by SHA-256 in its registry entry, together with the commit that produced it and the seeds it
used. The protocol shared by all arms is a single frozen object, so two arms differ only in the one
factor each experiment varies. Runs were executed on the Bocconi HPC. Supporting material: the
claim--evidence matrix, the negative-results record, the experiment registry and the
reproducibility notes in the project repository.}

\end{document}
